\makeatletter
\def\input@path{{./}{../}{../../}{../../../}{../../../../}{preamble/}{../preamble/}{../../preamble/}{../../../preamble/}}
\makeatother
\input{preamble_exams.tex}
% Only the shared exam preamble is needed; it already supplies TikZ and tables.
\tikzset{
 cNine survivor/.style={draw,align=center,font=\small,text width=65mm,minimum height=9mm},
 cNine selected/.style={cNine survivor,very thick,fill=black!8},
 cNine eliminated/.style={align=center,font=\small,text width=48mm},
 cNine keep/.style={->,thick},cNine discard/.style={->,dashed}}
\title{Grade 10 -- Term 1\\Decision-Making Under Uncertainty\\Exam Solutions -- Version 2C}
\author{}\date{}
\begin{document}
\maketitle
\setcounter{tocdepth}{1}\setcounter{secnumdepth}{2}
\phantomsection\label{toc}\tableofcontents
\section*{Learning objective}
Compare decision criteria and their limitations, and construct and interpret decision trees for staged filtering under uncertainty.
\section*{Criteria assessment}
This exam assesses C8 and C9, each worth 6 points. A \textbf{minimum of 4 points per criterion} is required to pass. Show all relevant calculations and justify your answers. Partial credit may be awarded for correct procedures despite arithmetic errors; answers without demonstrated procedures receive no credit. A 1-point penalty per criterion applies if the student's name is missing or the assessment is submitted late. C10 is assessed separately.
\begin{center}
\begin{tabular}{|c|c|c|}
				\hline
				\textbf{Question} & \textbf{C8} & \textbf{C9} \\ \hline
				1 & 3 & -- \\ \hline
				2 & 3 & -- \\ \hline
				3 & -- & 2 \\ \hline
				4 & -- & 2 \\ \hline
				5 & -- & 2 \\ \hline
				\textbf{Total} & \textbf{6} & \textbf{6} \\ \hline
			\end{tabular}
\end{center}
\section*{Evaluated criteria}
\begin{enumerate}[label=\textbf{C\arabic*},start=8]
\item Compares the Maximax, Maximin, Minimax Regret, and EMV/Bayes decision criteria by contrasting their information requirements, risk orientations, recommended alternatives, and limitations.
\item Constructs and interprets decision trees for staged filtering rules, applying decision criteria in the specified order to identify the final alternative.
\item Evaluates economic and financial decision-making scenarios by selecting and justifying the most appropriate decision criterion according to the information available, risk attitude, and decision context.
\end{enumerate}



\sectionScored{Camera Rental Plans}{C8-\sml{3}}{sec:q1}
A camera rental studio selects its inventory under uncertain demand. Profits are in thousand USD.

\vspace{-3mm}
\begin{center}\small\setlength{\tabcolsep}{4pt}
	\begin{tabular}{lrrr}\toprule
		Alternative & High & Normal & Low\\
		$P(S_j)$ & 0.2 & 0.5 & 0.3\\ \midrule
				Essential ($A_1$) & 115 & 100 & 75\\
		Shared ($A_2$) & 155 & 165 & 45\\
		Extensive ($A_3$) & 265 & 55 & -30\\
		Modular ($A_4$) & 190 & 145 & -10\\
		\bottomrule
	\end{tabular}
\end{center} \vspace*{-3mm}

Decision statistics (Regret = maximum regret):

\vspace{-3mm}
\begin{center}\small\setlength{\tabcolsep}{4pt}
	\begin{tabular}{lrrrr}\toprule
		Alt. & Max & Min & Regret & EMV\\ \midrule
				$A_1$ & 115 & 75 & 150 & 95.5\\
		$A_2$ & 165 & 45 & 110 & 127\\
		$A_3$ & 265 & -30 & 110 & 71.5\\
		$A_4$ & 190 & -10 & 85 & 107.5\\
		\bottomrule
	\end{tabular}
\end{center}
\vspace{-3mm}

\begin{enumerate}[label=\textbf{(\alph*)},leftmargin=*]
	\item Using the original camera rental payoffs, identify the alternative recommended by Maximax, Maximin, Minimax Regret, and EMV/Bayes, and report each criterion's deciding statistic value. \score{1}

	\item A subsidy $s$ (integer, $0\leq s\leq100$) increases the low-demand payoffs of $A_3$ and $A_4$ by $s$. An alternative is eligible only if all its payoffs are nonnegative. Determine each camera rental plan's eligibility based on the value of $s$. \score{0.5}

	\item Among eligible camera rental plans, find the smallest allowed subsidy $s$ at which the Maximin recommendation changes from $A_1$ to $A_4$ alone. Identify any tie at the threshold and distinguish it from a strict change. Respect the integer domain $0\leq s\leq100$; if no smallest value exists, state the threshold and the condition for a change. \score{1.5}
\end{enumerate}
\phantomsection\label{sol:q1a}
\subsection*{(a) Initial recommendations \score{1}}
\begin{center}
\begin{tabular}{lll}\toprule
Criterion & Recommended alternative(s) & Deciding statistic and value\\ \midrule
		Maximax & $A_3$ & Row maximum: 265\\
		Maximin & $A_1$ & Row minimum: 75\\
		Minimax Regret & $A_4$ & Maximum regret: 85\\
		EMV/Bayes & $A_2$ & EMV: 127\\
		\bottomrule\end{tabular}
\end{center}
All deciding values are in thousand USD. Maximax pursues the greatest upside; Maximin protects the worst outcome; Minimax Regret limits missed profit; EMV uses the forecast probabilities to maximize expected profit. Regret benchmarks use all four original alternatives, with state bests $(265,165,75)$; no eligibility restriction applies in (a).
\secInicio[sec:q1]


\phantomsection\label{sol:q1b}
\subsection*{(b) Subsidy and eligibility \score{0.5}}

The updated payoff table is:

\begin{center}\small\setlength{\tabcolsep}{4pt}
	\begin{tabular}{lrrr}\toprule
		Alternative & High & Normal & Low\\
		$P(S_j)$ & 0.2 & 0.5 & 0.3\\ \midrule
				Essential ($A_1$) & 115 & 100 & 75\\
		Shared ($A_2$) & 155 & 165 & 45\\
		Extensive ($A_3$) & 265 & 55 & $-30+s$\\
		Modular ($A_4$) & 190 & 145 & $-10+s$\\
		\bottomrule
	\end{tabular}
\end{center}

The updated decision statistics are:

\begin{center}\small\setlength{\tabcolsep}{4pt}
	\begin{tabular}{lrrrr}\toprule
		Alt. & Max & Min & Regret & EMV\\ \midrule
				$A_1$ & 115 & 75 & 150 & 95.5\\
		$A_2$ & 165 & 45 & 110 & 127\\
		$A_3$ & 265 & $\min(55,-30+s)$ & $\max(110,105-s)$ & $71.5+0.3s$\\
		$A_4$ & 190 & $-10+s$ & $\max(75,85-s)$ & $107.5+0.3s$\\
		\bottomrule
	\end{tabular}
\end{center}

An alternative is eligible only if all its payoffs are nonnegative.
For integer subsidies $0\le s\le100$:

\begin{center}
\begin{tabular}{lll}\toprule
Alternative & Eligibility condition & Eligible subsidies\\\midrule
$A_1$ & $75\ge0$ & $0\le s\le100$\\
$A_2$ & $45\ge0$ & $0\le s\le100$\\
$A_3$ & $-30+s\ge0$ & $30\le s\le100$\\
$A_4$ & $-10+s\ge0$ & $10\le s\le100$\\
\bottomrule
\end{tabular}
\end{center}

Thus, the eligible sets are:
\[
\begin{array}{ll}
0\le s\le9: & \{A_1,A_2\},\\
10\le s\le29: & \{A_4,A_1,A_2\},\\
30\le s\le100: & \{A_3,A_4,A_1,A_2\}.
\end{array}
\]

\secInicio[sec:q1]

\phantomsection\label{sol:q1c}
\subsection*{(c) Maximin change under subsidy $s$ \score{1.5}}

The minimum payoffs are
\[
m_1=75,\quad m_2=45,\quad
m_3(s)=\min(55,-30+s),\quad m_4(s)=-10+s.
\]

Since $m_3(s)\le55<75=m_1$ and $m_2=45<75$ throughout
the allowed range, only $A_4$ can challenge $A_1$.

For $A_4$ to beat $A_1$:
\[
-10+s>75
\quad\Longleftrightarrow\quad s>85.
\]

Thus, $A_1$ is the sole Maximin choice for $0\le s\le84$;
at $s=85$, $A_4$ and $A_1$ tie with minimum payoff 75;
and $A_4$ is the sole Maximin choice for $86\le s\le100$.
The tie at $s=85$ first changes the set of recommended plans;
a strict change to $A_4$ alone requires $s>85$.
The integer domain contains a smallest value satisfying this condition.
Therefore, the smallest allowed integer subsidy giving a strict change is
\[
\boxed{s=86}.
\]

At $s=86$, all four alternatives are eligible and their minimum
payoffs are $75,45,55,76$, respectively. Hence $A_4$ wins uniquely
with a guaranteed profit of 76 thousand USD.
\secInicio[sec:q1]



\sectionScored{Outdoor Cinema Refreshments}{C8-\sml{3}}{sec:q2}
A cinema manager selects a refreshment kiosk under uncertain audience size. Profits are in hundred USD.

\begin{center}\small\setlength{\tabcolsep}{4pt}
	\begin{tabular}{lrr}\toprule
		Alternative & High & Low\\
		$P(S_j)$ & 0.3 & 0.7\\ \midrule
				Portable ($C_1$) & 108 & 72\\
		Full-service ($C_2$) & 174 & 22\\
		Deluxe ($C_3$) & 204 & -8\\
		Fixed ($C_4$) & 158 & 54\\
		\bottomrule
	\end{tabular}
\end{center} \vspace{-3mm}

Decision statistics (Regret = maximum regret):

\begin{center}\small\setlength{\tabcolsep}{4pt}
	\begin{tabular}{lrrrr}\toprule
		Alternative & Max & Min & Regret & EMV\\ \midrule
				$C_1$ & 108 & 72 & 96 & 82.8\\
		$C_2$ & 174 & 22 & 50 & 67.6\\
		$C_3$ & 204 & -8 & 80 & 55.6\\
		$C_4$ & 158 & 54 & 46 & 85.2\\
		\bottomrule
	\end{tabular}
\end{center}

\begin{enumerate}[label=\textbf{(\alph*)},leftmargin=*]
	\item Using the original refreshment kiosk payoffs, identify the alternative recommended by Maximax, Maximin, Minimax Regret, and EMV/Bayes, and report each criterion's deciding statistic value. \score{1}

	\item A sponsor replaces $C_1$'s large-audience payoff of 108 with $g$ (integer, $108\leq g\leq180$). An alternative is eligible only if all its payoffs are at least 46. Determine each refreshment kiosk's eligibility based on the value of $g$. \score{0.5}

	\item Among eligible refreshment kiosks, find the smallest allowed sponsored payoff $g$ at which the Minimax Regret recommendation changes from $C_4$ to $C_1$ alone. Recalculate regret benchmarks using only eligible kiosks. Identify any tie at the threshold and distinguish it from a strict change. Respect the integer domain $108\leq g\leq180$; if no smallest value exists, state the threshold and the condition for a change. \score{1.5}
\end{enumerate}
\phantomsection\label{sol:q2a}
\subsection*{(a) Initial recommendations \score{1}}
\begin{center}
\begin{tabular}{lll}\toprule
Criterion & Recommended alternative(s) & Deciding statistic and value\\ \midrule
		Maximax & $C_3$ & Row maximum: 204\\
		Maximin & $C_1$ & Row minimum: 72\\
		Minimax Regret & $C_4$ & Maximum regret: 46\\
		EMV/Bayes & $C_4$ & EMV: 85.2\\
		\bottomrule\end{tabular}
\end{center}
All values are in hundred USD. Regret benchmarks use all four original kiosks, with state bests $(204,72)$; no eligibility restriction applies in (a). Only EMV uses the probabilities; the other criteria compare upside, downside, or missed profit.
\secInicio[sec:q2]



\phantomsection\label{sol:q2b}
\subsection*{(b) Sponsored payoff and eligibility \score{0.5}}

The updated payoff table is:

\begin{center}\small\setlength{\tabcolsep}{4pt}
	\begin{tabular}{lrr}\toprule
		Alternative & High & Low\\
		$P(S_j)$ & 0.3 & 0.7\\ \midrule
				Portable ($C_1$) & $g$ & 72\\
		Full-service ($C_2$) & 174 & 22\\
		Deluxe ($C_3$) & 204 & -8\\
		Fixed ($C_4$) & 158 & 54\\
		\bottomrule
	\end{tabular}
\end{center}

The updated decision statistics, before applying
eligibility restrictions, are:

\begin{center}\small\setlength{\tabcolsep}{4pt}
	\begin{tabular}{lrrrr}\toprule
		Alt. & Max & Min & Regret & EMV\\ \midrule
				$C_1$ & $g$ & 72 & $204-g$ & $0.3g+50.4$\\
		$C_2$ & 174 & 22 & 50 & 67.6\\
		$C_3$ & 204 & -8 & 80 & 55.6\\
		$C_4$ & 158 & 54 & 46 & 85.2\\
		\bottomrule
	\end{tabular}
\end{center}

An alternative is eligible only if both payoffs
are at least 46. For integers $108\le g\le180$:

\begin{center}
\begin{tabular}{lll}\toprule
Alternative & Eligibility condition & Eligible values\\\midrule
$C_1$ & $g\ge46,\ 72\ge46$ & $108\le g\le180$\\
$C_2$ & $22<46$ & None\\
$C_3$ & $-8<46$ & None\\
$C_4$ & $158\ge46,\ 54\ge46$ & $108\le g\le180$\\
\bottomrule
\end{tabular}
\end{center}

The eligible set is always $\{C_1,C_4\}$.
\secInicio[sec:q2]

\phantomsection\label{sol:q2c}
\subsection*{(c) Minimax Regret change under sponsored payoff $g$ \score{1.5}}

Recalculate regrets using only these alternatives,
with state benchmarks
\[
b_H=\max\{g,158\},\qquad b_L=72.
\]

Their maximum regrets are
\[
R(C_1)=\max\{158-g,0\},
\]
\[
R(C_4)=\max\{g-158,18\}.
\]

For $108\le g\le158$, these become
\[
R(C_1)=158-g,\qquad R(C_4)=18.
\]

For $C_1$ to be the sole Minimax Regret choice:
\[
158-g<18
\quad\Longleftrightarrow\quad g>140.
\]

At $g=140$, both alternatives have maximum regret 18.
Thus, $C_4$ is the sole Minimax Regret choice for $108\le g\le139$;
at $g=140$, the recommendation becomes the tie $\{C_1,C_4\}$;
and $C_1$ is the sole choice for $141\le g\le180$.
For $g>158$, $C_1$ has zero regret and $C_4$ has regret at least 18,
so no further change occurs.
The tie first changes the set of recommended kiosks, but a strict
change to $C_1$ alone requires $g>140$.
The integer domain contains a smallest value satisfying this condition.
Therefore, the smallest allowed integer giving a strict change is
\[
\boxed{g=141}.
\]

At this value, the eligible payoff table is:

\begin{center}\small\setlength{\tabcolsep}{4pt}
\begin{tabular}{lrr}\toprule
Alternative & High & Low\\ \midrule
$C_1$ & 141 & 72\\
$C_4$ & 158 & 54\\
\bottomrule
\end{tabular}
\end{center}

The corresponding maximum regrets are
\[
R(C_1)=17,\qquad R(C_4)=18.
\]

Thus, $C_1$ wins uniquely with $g=141$,
corresponding to \$14,100 in large-audience profit.
\secInicio[sec:q2]



\sectionScored{Garden Centre Stock}{C9-\sml{2}}{sec:q3}
A garden centre selects a seedling stock plan $D_1$--$D_6$. Net profits are in hundred USD. The decision statistics are summarized below:

\vspace{-3mm}
\begin{center}\small\setlength{\tabcolsep}{4pt}
	\begin{tabular}{lrrrr}\toprule
		Alternative & Max & Min & Regret & EMV\\ \midrule
		$D_1$ & 164 & 68 & 72 & 137.6\\
		$D_2$ & 188 & 44 & 48 & 154.4\\
		$D_3$ & 140 & 92 & 96 & 123.2\\
		$D_4$ & 212 & 20 & 72 & 154.4\\
		$D_5$ & 116 & 68 & 120 & 99.2\\
		$D_6$ & 236 & -4 & 96 & 154.4\\
		\bottomrule
	\end{tabular}
\end{center}
\vspace{-3mm}

Apply these stages in order:
\begin{enumerate}[leftmargin=*]
	\item EMV/Bayes: Best 3.
	\item Maximax: Best 2.
	\item Maximin: Final selection.
\end{enumerate}
Illustrate each filtering stage, showing survivors and eliminations.
\phantomsection\label{sol:q3table}
\subsection*{Supporting filtering table}
\begin{center}\small
\begin{tabularx}{\textwidth}{@{}c >{\raggedright\arraybackslash}p{0.24\textwidth} >{\raggedright\arraybackslash}X >{\raggedright\arraybackslash}p{0.23\textwidth}@{}}
\toprule
Step & Rule & Survivors (relevant values) & Eliminated (values)\\\midrule
0 & Initial alternatives & $D_6,D_4,D_2,D_1,D_3,D_5$ & ---\\
1 & EMV/Bayes: Best 3 & $D_6(154.4)$, $D_4(154.4)$, $D_2(154.4)$ & $D_1(137.6)$, $D_3(123.2)$, $D_5(99.2)$\\
2 & Maximax: Best 2 & $D_6(236)$, $D_4(212)$ & $D_2(188)$\\
3 & Maximin: Final selection & \textbf{Selected:} $D_4(20)$ & $D_6(-4)$\\
\bottomrule\end{tabularx}
\end{center}
The highest EMV, 154.4, is shared by $D_6,D_4,D_2$; all three fill the Best 3 places. Maximax retains $D_6,D_4$ with maxima 236 and 212, eliminating $D_2$ with 188. Maximin selects $D_4$ because $20>-4$. There are no cutoff or final ties beyond the shared first-stage EMV.
\secInicio[sec:q3]
\clearpage
\phantomsection\label{sol:q3tree}
\subsection*{Filtering tree and interpretation}
\begin{center}
\begin{tikzpicture}[>=Stealth]
\node[cNine survivor] (s0) at (0,0) {All six $D_i$};
\node[cNine survivor] (s1) at (0,-1.65) {$D_6,D_4,D_2$};
\draw[cNine keep] (s0) -- node[right,font=\small] {EMV/Bayes: Best 3} (s1);
\node[cNine eliminated] (e1) at (7.5,-1.65) {Eliminate $D_1,D_3,D_5$};
\draw[cNine discard] (s1) -- (e1);
\node[cNine survivor] (s2) at (0,-3.3) {$D_6,D_4$};
\draw[cNine keep] (s1) -- node[right,font=\small] {Maximax: Best 2} (s2);
\node[cNine eliminated] (e2) at (7.5,-3.3) {Eliminate $D_2$};
\draw[cNine discard] (s2) -- (e2);
\node[cNine selected] (s3) at (0,-4.95) {\textbf{Selected:} $D_4$};
\draw[cNine keep] (s2) -- node[right,font=\small] {Maximin: Final selection} (s3);
\node[cNine eliminated] (e3) at (7.5,-4.95) {Eliminate $D_6$};
\draw[cNine discard] (s3) -- (e3);
\end{tikzpicture}
\end{center}
The garden centre selects $D_4$: it survives the EMV and Maximax filters, then its minimum profit of 20 hundred USD exceeds $D_6$'s -4.
\secInicio[sec:q3]

\sectionScored{Climbing Gym Expansion}{C9-\sml{2}}{sec:q4}
A climbing gym selects an expansion plan $F_1$--$F_6$. Net profits are in thousand USD. The decision statistics are summarized below:
\vspace{-3mm}
\begin{center}\small\setlength{\tabcolsep}{4pt}
	\begin{tabular}{lrrrr}\toprule
		Alternative & Max & Min & Regret & EMV\\ \midrule
		$F_1$ & 213 & 59 & 55 & 154.7\\
		$F_2$ & 169 & 103 & 99 & 144.8\\
		$F_3$ & 191 & 81 & 77 & 164.6\\
		$F_4$ & 268 & 4 & 99 & 138.2\\
		$F_5$ & 147 & 92 & 121 & 125\\
		$F_6$ & 235 & 37 & 66 & 145.9\\
		\bottomrule
	\end{tabular}
\end{center}
\vspace{-3mm}
Apply these stages in order:
\begin{enumerate}[leftmargin=*]
	\item Minimax Regret: Best 5.
	\item EMV/Bayes: Best 4.
	\item Maximax: Best 2.
	\item Maximin: Final selection.
\end{enumerate}
Illustrate each filtering stage, showing survivors and eliminations.
\phantomsection\label{sol:q4table}
\subsection*{Supporting filtering table}
\begin{center}\small
\begin{tabularx}{\textwidth}{@{}c >{\raggedright\arraybackslash}p{0.24\textwidth} >{\raggedright\arraybackslash}X >{\raggedright\arraybackslash}p{0.23\textwidth}@{}}
\toprule
Step & Rule & Survivors (relevant values) & Eliminated (values)\\\midrule
0 & Initial alternatives & $F_4,F_6,F_1,F_3,F_2,F_5$ & ---\\
1 & Minimax Regret: keep 5 lowest & $F_1(55)$, $F_6(66)$, $F_3(77)$, $F_4(99)$, $F_2(99)$ & $F_5(121)$\\
2 & EMV/Bayes: keep 4 highest & $F_3(164.6)$, $F_1(154.7)$, $F_6(145.9)$, $F_2(144.8)$ & $F_4(138.2)$\\
3 & Maximax: keep 2 highest Max & $F_6(235)$, $F_1(213)$ & $F_3(191)$, $F_2(169)$\\
4 & Maximin: select highest Min & \textbf{Selected:} $F_1(59)$ & $F_6(37)$\\
\bottomrule\end{tabularx}
\end{center}
At the first cutoff, $F_4$ and $F_2$ both have regret 99, and both survive. Later cutoffs have no ties.
\secInicio[sec:q4]

\phantomsection\label{sol:q4tree}
\subsection*{Filtering tree and interpretation}
\begin{center}
\begin{tikzpicture}[>=Stealth]
\node[cNine survivor] (s0) at (0,0) {All six $F_i$};
\node[cNine survivor] (s1) at (0,-1.65) {$F_1,F_6,F_3,F_4,F_2$};
\draw[cNine keep] (s0) -- node[right,font=\small] {Minimax Regret: keep 5 lowest} (s1);
\node[cNine eliminated] (e1) at (7.5,-1.65) {Eliminate $F_5$};
\draw[cNine discard] (s1) -- (e1);
\node[cNine survivor] (s2) at (0,-3.3) {$F_3,F_1,F_6,F_2$};
\draw[cNine keep] (s1) -- node[right,font=\small] {EMV/Bayes: keep 4 highest} (s2);
\node[cNine eliminated] (e2) at (7.5,-3.3) {Eliminate $F_4$};
\draw[cNine discard] (s2) -- (e2);
\node[cNine survivor] (s3) at (0,-4.949999999999999) {$F_6,F_1$};
\draw[cNine keep] (s2) -- node[right,font=\small] {Maximax: keep 2 highest Max} (s3);
\node[cNine eliminated] (e3) at (7.5,-4.949999999999999) {Eliminate $F_3,F_2$};
\draw[cNine discard] (s3) -- (e3);
\node[cNine selected] (s4) at (0,-6.6) {\textbf{Selected:} $F_1$};
\draw[cNine keep] (s3) -- node[right,font=\small] {Maximin: select highest Min} (s4);
\node[cNine eliminated] (e4) at (7.5,-6.6) {Eliminate $F_6$};
\draw[cNine discard] (s4) -- (e4);
\end{tikzpicture}
\end{center}
The climbing gym selects $F_1$: after the regret, EMV, and upside filters, its minimum profit of 59 thousand USD exceeds $F_6$'s 37.
\secInicio[sec:q4]
\clearpage
\sectionScored{Art Materials Orders}{C9-\sml{2}}{sec:q5}
An art materials distributor selects an order plan $H_1$--$H_6$. Net profits are in hundred USD. The decision statistics are summarized below:
\vspace{-3mm}
\begin{center}\small\setlength{\tabcolsep}{4pt}
	\begin{tabular}{lrrrr}\toprule
		Alternative & Max & Min & Regret & EMV\\ \midrule
		$H_1$ & 129 & 81 & 72 & 111.4\\
		$H_2$ & 145 & 65 & 56 & 124.2\\
		$H_3$ & 177 & 33 & 48 & 112.2\\
		$H_4$ & 113 & 73 & 88 & 97\\
		$H_5$ & 201 & 9 & 72 & 106.6\\
		$H_6$ & 161 & 49 & 40 & 118.6\\
		\bottomrule
	\end{tabular}
\end{center}
\vspace{-3mm}
Apply these stages in order:
\begin{enumerate}[leftmargin=*]
	\item Maximin: Best 4.
	\item Minimax Regret: Best 3.
	\item Maximax: Best 2.
	\item EMV/Bayes: Final selection.
\end{enumerate}
Illustrate each filtering stage, showing survivors and eliminations.
\phantomsection\label{sol:q5table}
\subsection*{Supporting filtering table}
\begin{center}\small
\begin{tabularx}{\textwidth}{@{}c >{\raggedright\arraybackslash}p{0.24\textwidth} >{\raggedright\arraybackslash}X >{\raggedright\arraybackslash}p{0.23\textwidth}@{}}
\toprule
Step & Rule & Survivors (relevant values) & Eliminated (values)\\\midrule
0 & Initial alternatives & $H_5,H_3,H_6,H_2,H_1,H_4$ & ---\\
1 & Maximin: keep 4 highest Min & $H_1(81)$, $H_4(73)$, $H_2(65)$, $H_6(49)$ & $H_3(33)$, $H_5(9)$\\
2 & Minimax Regret: keep 3 lowest & $H_6(40)$, $H_2(56)$, $H_1(72)$ & $H_4(88)$\\
3 & Maximax: keep 2 highest Max & $H_6(161)$, $H_2(145)$ & $H_1(129)$\\
4 & EMV/Bayes: select highest EMV & \textbf{Selected:} $H_2(124.2)$ & $H_6(118.6)$\\
\bottomrule\end{tabularx}
\end{center}
The cutoffs are Min 49, maximum regret 72, and Max 145; the final EMV comparison is 124.2 versus 118.6. There are no cutoff or final ties.
\secInicio[sec:q5]

\phantomsection\label{sol:q5tree}
\subsection*{Filtering tree and interpretation}
\begin{center}
\begin{tikzpicture}[>=Stealth]
\node[cNine survivor] (s0) at (0,0) {All six $H_i$};
\node[cNine survivor] (s1) at (0,-1.65) {$H_1,H_4,H_2,H_6$};
\draw[cNine keep] (s0) -- node[right,font=\small] {Maximin: keep 4 highest Min} (s1);
\node[cNine eliminated] (e1) at (7.5,-1.65) {Eliminate $H_3,H_5$};
\draw[cNine discard] (s1) -- (e1);
\node[cNine survivor] (s2) at (0,-3.3) {$H_6,H_2,H_1$};
\draw[cNine keep] (s1) -- node[right,font=\small] {Minimax Regret: keep 3 lowest} (s2);
\node[cNine eliminated] (e2) at (7.5,-3.3) {Eliminate $H_4$};
\draw[cNine discard] (s2) -- (e2);
\node[cNine survivor] (s3) at (0,-4.949999999999999) {$H_6,H_2$};
\draw[cNine keep] (s2) -- node[right,font=\small] {Maximax: keep 2 highest Max} (s3);
\node[cNine eliminated] (e3) at (7.5,-4.949999999999999) {Eliminate $H_1$};
\draw[cNine discard] (s3) -- (e3);
\node[cNine selected] (s4) at (0,-6.6) {\textbf{Selected:} $H_2$};
\draw[cNine keep] (s3) -- node[right,font=\small] {EMV/Bayes: select highest EMV} (s4);
\node[cNine eliminated] (e4) at (7.5,-6.6) {Eliminate $H_6$};
\draw[cNine discard] (s4) -- (e4);
\end{tikzpicture}
\end{center}
The art materials distributor selects $H_2$: after the downside, regret, and upside filters, its EMV of 124.2 hundred USD exceeds $H_6$'s 118.6.
\secInicio[sec:q5]
\section*{C10: Separate activity reminder (unscored)}
Complete the individual \textbf{Interactive Decision Webpage: C10 Practice} activity and submit the deployed webpage URL. See the separate activity for instructions.

\href{https://nolopez338.github.io/global_economics/grade10/src/pages/term1/g10-t1-c10dashboard.html}{C10 example solution.}
\end{document}
